% Current verified additions, 2026-11-08.
% Insert inside the manuscript's thebibliography environment.

\bibitem{UeyamaIsolated}
K. Ueyama,
\emph{Noncommutative graded Gorenstein isolated singularities},
in \emph{Proceedings of the 45th Symposium on Ring Theory and Representation Theory}
(Shinshu University, September 7--9, 2012),
K. Koike (ed.), February 2013, 158--163.
Theorem 7, printed p.~161.
\url{https://ring-theory-japan.com/ring/oldmeeting/2012/report2012/proceedings.pdf}.

\bibitem{ZhangTwist}
J.~J. Zhang,
\emph{Twisted graded algebras and equivalences of graded categories},
Proceedings of the London Mathematical Society (3) \textbf{72} (1996), no.~2, 281--311.
Theorems 1.1 and 1.3.
\url{https://doi.org/10.1112/plms/s3-72.2.281}.
\url{https://deepblue.lib.umich.edu/bitstream/handle/2027.42/135651/plms0281.pdf}.

\bibitem{RRZ}
M. Reyes, D. Rogalski and J.~J. Zhang,
\emph{Skew Calabi--Yau algebras and homological identities},
Advances in Mathematics \textbf{264} (2014), 308--354.
Theorems 0.3 and 5.4; Example 5.5.
\url{https://doi.org/10.1016/j.aim.2014.07.010}.
\url{https://arxiv.org/abs/1302.0437}.

\bibitem{Griffiths}
P.~A. Griffiths,
\emph{On the periods of certain rational integrals.~I},
Annals of Mathematics (2) \textbf{90} (1969), no.~3, 460--495.
Section 8, especially Theorems 8.1 and 8.3 and formula (8.6).
\url{https://doi.org/10.2307/1970746}.
\url{https://www.cimat.mx/~luis/seminarios/Teoria-de-Hodge/period_integral_I.pdf}.

\bibitem{KuznetsovFCY}
A. Kuznetsov,
\emph{Calabi--Yau and fractional Calabi--Yau categories},
Journal f\"ur die reine und angewandte Mathematik \textbf{753} (2019), 239--267.
Lemmas 2.6, 3.13 and 3.19; Corollary 4.1.
\url{https://doi.org/10.1515/crelle-2017-0004}.

\bibitem{CHL}
C.-H. Cho, H. Hong and S.-C. Lau,
\emph{Noncommutative homological mirror functor},
Memoirs of the American Mathematical Society \textbf{271} (2021), no.~1326.
Chapter 5, especially Theorems 5.15 and 5.16.
\url{https://doi.org/10.1090/memo/1326}.
\url{https://arxiv.org/abs/1512.07128}.

\bibitem{Lee}
S. Lee,
\emph{Noncommutative homological mirror symmetry of elliptic curves},
Kyoto Journal of Mathematics \textbf{61} (2021), no.~3, 723--743.
Theorem 6.3 and the corollary in the introduction.
\url{https://arxiv.org/abs/2004.09145}.

\bibitem{CisinskiTabuada}
D.-C. Cisinski and G. Tabuada,
\emph{Lefschetz and Hirzebruch--Riemann--Roch formulas via noncommutative motives},
Journal of Noncommutative Geometry \textbf{8} (2014), 1171--1190.
Section 4 and Example 5.6.
\url{https://doi.org/10.4171/JNCG/183}.
\url{https://www.math.univ-toulouse.fr/~dcisinsk/ncLHRR.pdf}.

\bibitem{MathewHH}
A. Mathew,
\emph{Kaledin's degeneration theorem and topological Hochschild homology},
Geometry \& Topology \textbf{24} (2020), no.~6, 2675--2708.
Theorem 1.3.
\url{https://doi.org/10.2140/gt.2020.24.2675}.
\url{https://arxiv.org/abs/1710.09045}.

\bibitem{PerryIHC}
A. Perry,
\emph{The integral Hodge conjecture for two-dimensional Calabi--Yau categories},
Compositio Mathematica \textbf{158} (2022), no.~2, 287--333.
Lemmas 3.4 and 3.6; Theorem 3.5.
\url{https://doi.org/10.1112/S0010437X22007266}.

\bibitem{FuHwangCI}
B. Fu and J.-M. Hwang,
\emph{Isotrivial VMRT-structures of complete intersection type},
Asian Journal of Mathematics \textbf{22} (2018), 333--354.
Lemma 4.1 (Br\"uckmannの消滅定理を引用).
\url{https://doi.org/10.4310/AJM.2018.v22.n2.a10}.
\bibitem{BorelRemmert}
A. Borel and R. Remmert,
\emph{\"Uber kompakte homogene K\"ahlersche Mannigfaltigkeiten},
Mathematische Annalen \textbf{145} (1962), 429--439, Theorem I.
\url{https://wglobke.github.io/publications/BorelRemmert1962.pdf}.

\bibitem{DeLaetLeBruyn}
K. De Laet and L. Le Bruyn,
\emph{The geometry of representations of 3-dimensional Sklyanin algebras},
Algebras and Representation Theory \textbf{18} (2015), 761--776.
arXiv:1405.1158v1, Theorem 1 and Section 2 (canonical central cubic formula);
the result is a reformulation of the Artin--Schelter and Smith--Tate formula.
\url{https://arxiv.org/abs/1405.1158}.

\bibitem{ArtebaniDolgachev}
M. Artebani and I. Dolgachev,
\emph{The Hesse pencil of plane cubic curves},
L'Enseignement Math\'ematique (2) \textbf{55} (2009), no.~3--4, 235--273,
Lemma 1 (every smooth plane cubic admits a Hesse canonical form).
\url{https://doi.org/10.4171/LEM/55-3-3};
\url{https://arxiv.org/abs/math/0611590}.
